\section{The transition at small tolerance}\label{sec:limit}
Write $a_h(\delta)=m_h/N_h$ and let $H_*(B,\delta)$ be the first $H\ge h_0$ with $C_{p,\delta}(h_0,H)>B$. This index bounds the \emph{exponent}: the order of the root is at most $p^{H_*}$. The comparison is strict because the extremum in Theorem~\ref{thm:exact-obstruction} is attained. The formula is readily computable by binary search for $m_h$ with integer or rational arithmetic.

\paragraph{A finite example.}
At $p=2$, $h_0=1$ and $\delta=1/4$ the optimal costs are as follows.
\begin{center}\begin{tabular}{@{}lrrrrrr@{}}\toprule
Shell exponent $h$ & 1 & 2 & 3 & 4 & 5 & 6\\
\midrule
Primitive roots $N_h$ & 1&2&4&8&16&32\\
Minimum marks $m_h$ &1&1&2&2&3&4\\
Cumulative budget & $1$ & $3/2$ & $2$ & $9/4$ & $39/16$ & $41/16$\\
\bottomrule\end{tabular}\end{center}
A budget of exactly $2$ can block orders $2,4,8$. Every packet of that budget admits a successful point of order at most $16$. At order $8$, one mark attains minimum loss exactly $1/4$ across the primitive roots; the definition counts equality as success, so blocking this shell requires a second mark.

\begin{theorem}[The fixed-budget transition]\label{thm:transition}
Set $s_h=p^{1-h}$ and
\[
 B_0=\sum_{h=h_0}^\infty s_h=\frac{p^{2-h_0}}{p-1}.
\]
For fixed $p,h_0,B\ge0$ (with $p$ prime and $h_0\ge1$), the least guaranteed exponent as $\delta\downarrow0$ obeys the following three laws.
\begin{enumerate}
\item If $B<B_0$, it eventually equals the first $H$ with $\sum_{h=h_0}^H s_h>B$.
\item If $B=B_0$, put $t=\lceil1+\log_p(1/\delta)\rceil$. For all sufficiently small positive $\delta$, the answer is $t$ for odd $p$ and $t+1$ for $p=2$.
\item If $B>B_0$, then $\displaystyle\lim_{\delta\downarrow0}\delta\log H_*(B,\delta)=B-B_0$.
\end{enumerate}
\end{theorem}
\begin{proof}
For each fixed shell, $a_h(\delta)=s_h$ for every sufficiently small $\delta\ge0$. Indeed for $h\ge2$ the first nonzero finite score occurs at $L=p-1$, where $S_h(p-1)=s_h$. For $h=1$ all leaves must be marked. This proves the subcritical case and shows that $H_*(B,\delta)\to\infty$ when $B\ge B_0$.

The estimates in Corollary~\ref{cor:exponential-depth} give the uniform summable bound, for $h$ sufficiently large and $0\le\delta\le1$,
\[
 0\le a_h(\delta)-\delta/h
 \ll_p \frac{\log(2h)}{h^2}.
\]
For clarity, the upper bound uses $q=h-\lceil3\log_p h\rceil$ and
$a_h\le \delta/q +p^q/(qN_h)+1/N_h$. The lower bound follows from $S_h(L)\le(h-1)L/N_h$ for $L<N_h$, and also holds if $m_h=N_h$. Dominated convergence on the positive integers now yields, along \emph{any} $H\to\infty$, $\delta\downarrow0$,
\[
 C_{p,\delta}(h_0,H)-\delta\log H\longrightarrow B_0.
\]
Here $\sum_{h=h_0}^H 1/h-\log H$ is bounded. At $H=H_*$ the overshoot above $B$ is at most $a_H(\delta)$ once $H>h_0$. The uniform bound shows that this tends to zero, proving the supercritical case.

For the critical case take $\delta$ small enough that $t>\max(2,h_0)$. All earlier costs equal $s_h$. At $h=t$ the tolerance is in $[s_t,p s_t)$. All $N_j$ are multiples of $p-1$, so the score is constant between consecutive multiples of $p-1$ in the mark count. Hence $a_t\ge2s_t$. The remaining baseline tail is $\sum_{h>t}s_h=s_t/(p-1)$. For $p>2$ the extra cost already exceeds this tail and $H_*=t$. For $p=2$, the exact scores at $L=1,2$ are $s_t,4s_t$ since $t\ge3$. Thus $a_t=2s_t$, the cumulative cost is exactly $B_0$, and one further positive-cost shell makes it strictly larger.
\end{proof}

For $p=2$ and $h_0=1$, the critical budget is $2$. Raising the budget from exactly $2$ to $3$ changes the asymptotic exponent from $\lceil1+\log_2(1/\delta)\rceil+1$ to $\exp((1+o(1))/\delta)$. Thus the failure of a packet-size-independent polynomial bound already occurs with \emph{fixed} budget and arbitrarily small tolerance.
